% Copyright (c) 2021-2027 Patrick Hall

\documentclass[fleqn]{article}
\renewcommand\refname{}
\title{Responsible Machine Learning\\\Large{Assignment 2}\\\Large{10 points}}

\usepackage{graphicx}
\usepackage{fullpage}
\usepackage{pdfpages}
\usepackage{amsmath}
\usepackage{amssymb}
\usepackage{mathtools}
\usepackage{MnSymbol}
\usepackage{enumerate}
\usepackage{setspace}
\usepackage[colorlinks, breaklinks=true]{hyperref} 
\usepackage{float}
\usepackage{caption}
\usepackage{subcaption}
\usepackage{multicol}
\usepackage{color}
\usepackage{listings}
\usepackage{csvsimple}
\usepackage{algorithm}
\usepackage{algorithmic}
\usepackage{verbatim}
\usepackage{mdframed}
\usepackage{changepage}
\usepackage[top=1in, bottom=1in, left=1in, right=1in]{geometry}

\begin{document}

\maketitle

\noindent In Assignment 2, analyze and explain machine learning (ML) models developed for Assignment 1. You may work alone or in a group. Submit your work in a public GitHub repository, or in a private repository shared with the instructor (\texttt{jphall663}). You may use AI tools, but you are responsible for validating and understanding your results.\\

\noindent For each section, consider whether the explanatory results make sense given domain knowledge and whether differences between the models are logical and informative.\\

\section{Calculate and Plot Global Feature Importance.}\label{global_fi}

Use regression coefficients, Shapley values, or another technically sound method to calculate global feature importance for your models. Plot the results as bar charts and compare how the models rank and weight input features.\\

\section{Calculate and Plot Local Feature Importance.}\label{local_fi}

Calculate local feature importance for individuals at the 10\textsuperscript{th}, 50\textsuperscript{th}, and 90\textsuperscript{th} percentiles of predicted probability for each model.\\

Plot the local feature importance values and compare how the models explain individuals at similar points in their predicted probability distributions.

\section{Calculate and Plot Feature Behavior.}

Plot partial dependence for each main effect feature under each model. You may also plot interaction contours or surfaces, such as interactions identified by EBM. Alternative techniques include accumulated local effects (ALE), individual conditional expectation (ICE), and EBM feature scores.

\section{Submit Code Results.}

Submit your explanatory analysis and results in your GitHub repository. Report holdout performance for each model, rank the models using an appropriate holdout metric, and briefly explain the ranking.\\

\noindent Global feature importance plots: 2 points.\\
\noindent Local feature importance plots: 2 points.\\
\noindent Feature behavior plots: 2 points.\\
\noindent Ranking on holdout data: 4 points.\\
\noindent Total: 10 points.\\

\noindent \textbf{Deadline: TBD.}\\

\end{document}

